\documentclass[11pt,a4paper]{article}
\usepackage[T1]{fontenc}
\usepackage[utf8]{inputenc}
\usepackage{lmodern}
\usepackage[margin=2.5cm]{geometry}
\usepackage{booktabs}
\usepackage{tabularx}
\usepackage{longtable}
\usepackage{enumitem}
\usepackage{listings}
\usepackage{xcolor}
\usepackage[hidelinks]{hyperref}

\lstset{
  basicstyle=\ttfamily\small,
  columns=fullflexible,
  keepspaces=true,
  frame=single,
  rulecolor=\color{gray!50},
  backgroundcolor=\color{gray!8},
  breaklines=true,
  upquote=true,
  showstringspaces=false,
  extendedchars=true,
  inputencoding=utf8,
  literate={é}{{\'e}}1 {è}{{\`e}}1 {à}{{\`a}}1 {ç}{{\c c}}1 {É}{{\'E}}1,
  xleftmargin=0.5em,
  xrightmargin=0.5em
}
\setlength{\emergencystretch}{3em}
\newcommand{\kbd}[1]{\texttt{#1}}
\newcolumntype{L}{>{\raggedright\arraybackslash}X}

\title{LeRobot (Python, English names) User Manual}
\author{}
\date{10 November 2026}

\begin{document}
\maketitle

LeRobot is a drawing robot you steer on screen: either live, with the mouse
and keyboard (Draw), or by writing a short Python program of robot commands
(such as \kbd{my\_drawing.py}). This is the edition with English command
names, used by the book \emph{First Programming Lessons in Python}: the
robot moves forward with \kbd{fd()}, turns a quarter turn to the right with
\kbd{qtr()}, puts its pen down with \kbd{pd()}, and so on. It is the same
robot as the French edition (folder \kbd{LeRobot\_Python\_FR}, commands
\kbd{av()}, \kbd{pqtd()}, \kbd{bc()}, \dots), itself translated from the
OCaml and 1998 Caml Light versions: same window, same drawings. It lives in
the folder \kbd{LeRobot\_Python\_EN} of the \kbd{LeRobot} folder, which can
be anywhere (commands below are typed in a Terminal window opened in the
\kbd{LeRobot} folder); the installation is described in the \emph{LeRobot
(Python, English names) Installation Manual} (\kbd{LeRobot\_Install.pdf}).

\tableofcontents

%======================================================================
\section{Getting started}

You need a Python linked to Tk 8.6 or later, for example the one from
python.org (see the installation manual). There is nothing to compile and
no XQuartz: the robot window is an ordinary macOS window.

There are three ways to start a program:
\begin{itemize}
\item \textbf{The application} \kbd{LeRobotPythonEN.app} (in
  \kbd{LeRobot\_Python\_EN}): double-click it to run Draw; drop a \kbd{.py}
  file on its icon to run that program.
\item \textbf{The script} \kbd{robot.sh}, in a Terminal window opened in
  \kbd{LeRobot\_Python\_EN}: \kbd{./robot.sh path/to/Program.py}. The program
  can be in any folder.
\item \textbf{Python directly}, for the files of the \kbd{LeRobot\_Python\_EN}
  folder itself: \kbd{python3 Draw.py}, \kbd{python3 my\_drawing.py}.
\end{itemize}
If the \kbd{python3} of the Terminal is the one of macOS (linked to the old
Tk 8.5), LeRobot notices it and starts itself again with the python.org
Python; nothing has to be done.

The window is $800 \times 700$ points. The large white area is where the
robot draws; the strip along the bottom is the message bar.

%======================================================================
\section{Part 1 --- Interactive drawing (Draw)}

In Draw you give the robot commands with single keys and move it with
the mouse or the arrow keys; press Return when you are done, then \kbd{y} to
save the drawing as a program.

\subsection{Launch}
Double-click \kbd{LeRobotPythonEN.app}, or:
\begin{lstlisting}[language=bash]
cd LeRobot_Python_EN && python3 Draw.py
\end{lstlisting}
The robot appears in the centre, pointing north. The window comes to the
front with the keyboard focus; if keys do nothing, click in the window.

\subsection{The window}
\begin{itemize}
\item \textbf{Drawing area} --- the large white area.
\item \textbf{Message bar} (bottom) --- a reminder of the keys, ending with
  \kbd{* :pc w :paint\ \ Return :end}.
\item \textbf{Colour box} (small robot near the right of the bar) --- shows
  the current pen colour.
\item \textbf{Return icon} (far right) --- the Return symbol; it blinks when
  the robot waits for a key.
\end{itemize}

\subsection{Moving the robot with the mouse or the arrow keys}
Click on a grid point (or move the pointer with the button held down) and
the robot jumps to the nearest grid point, without drawing. The arrow keys
move it one grid step left, right, up or down, also without drawing. When
the pointer moves with the button released, the robot does not move.

If the robot would go past the edge of the drawing area, it stays where it
is and the terminal prints \kbd{** EXIT FORBIDDEN ** translate}, with a
beep-beep. This is only a warning.

\subsection{Keyboard commands}
The robot draws only when its pen is down (\kbd{d}; it is down at the
start). One step is one grid square.

\begin{center}
\begin{tabularx}{\linewidth}{@{}l l L@{}}
\toprule
\textbf{Key} & \textbf{Command} & \textbf{What it does} \\
\midrule
\kbd{d} & pd & \underline{P}en \underline{d}own: the robot draws as it moves \\
\kbd{u} & pu & \underline{P}en \underline{u}p: the robot moves without drawing \\
\kbd{f} & fd & Move \underline{f}orwar\underline{d} one step \\
\kbd{)} & cr & \underline{C}urve to the \underline{r}ight (a quarter circle) \\
\kbd{(} & cl & \underline{C}urve to the \underline{l}eft (a quarter circle) \\
\kbd{+} & qtr & \underline{Q}uarter \underline{t}urn \underline{r}ight ($90^\circ$) \\
\kbd{-} & qtl & \underline{Q}uarter \underline{t}urn \underline{l}eft ($90^\circ$) \\
\kbd{r} & etr & \underline{E}ighth \underline{t}urn \underline{r}ight ($45^\circ$) \\
\kbd{l} & etl & \underline{E}ighth \underline{t}urn \underline{l}eft ($45^\circ$) \\
\kbd{c} & dc & \underline{D}raw a small \underline{c}ross (+) \\
\kbd{p} & dp & \underline{D}raw a \underline{p}oint \\
\kbd{.} & dg & \underline{D}raw the \underline{g}rid points \\
\kbd{*} & pc & Next \underline{p}en \underline{c}olour (shown in the colour box) \\
\kbd{w} & paint & Paint the closed area behind the robot with the pen colour \\
Arrow keys & translate & Move one step left, right, up or down, without drawing \\
Backspace & --- & Undo the last command \\
Return & --- & Finish the drawing \\
\bottomrule
\end{tabularx}
\end{center}

Colours cycle in this order: white, yellow, blue, violet, green, red,
indigo, black. To paint, put the robot so that its back is inside a closed
area, then press \kbd{w}.

\subsection{Finishing and saving}
\begin{enumerate}
\item Press \textbf{Return}. The bar asks:
  \emph{Do you want to save in my\_drawing.py? (y/n)}
\item Press \kbd{y} to save. Every command you used, including mouse and
  arrow moves (written as \kbd{translate(dx, dy)}) and colour changes
  (written as \kbd{pc(n)} with the colour name in a comment), is written to
  \kbd{my\_drawing.py} as a Python program (see Part~2), and the window
  closes.
\item Press any other key to close without saving.
\end{enumerate}
Saving always overwrites \kbd{my\_drawing.py}, in the folder of
Draw. To keep an earlier drawing, copy that file under another name
first.

%======================================================================
\section{Part 2 --- Running a robot program}

A robot program is a Python file that lists robot commands in order. It
draws the picture by itself, then waits so you can replay it at another
speed or close it with Esc.

\subsection{Run my\_drawing.py}
\begin{lstlisting}[language=bash]
cd LeRobot_Python_EN && python3 my_drawing.py
\end{lstlisting}
or \kbd{make my\_drawing}, or drop \kbd{my\_drawing.py} on
\kbd{LeRobotPythonEN.app}. A saved drawing looks like this:
\begin{lstlisting}[language=Python]
# My lovely drawing
from LeRobot import *

init_pict(); gc()
fd(); fd(); qtr(); fd(); fd(); pc(2)  # blue
fd(); translate(5, 3); fd(); cl(); fd(); st()
\end{lstlisting}

\subsection{Write your own program}
Create a file anywhere, for example \kbd{house.py}:
\begin{lstlisting}[language=Python]
from LibRobot import *

# Draw a square
fd(); fd(); qtr(); fd(); fd(); qtr()
fd(); fd(); qtr(); fd(); fd(); qtr()
st()
\end{lstlisting}
Run it with \kbd{./robot.sh path/to/house.py} from the
\kbd{LeRobot\_Python\_EN} folder, or drop it on \kbd{LeRobotPythonEN.app}.

Every program follows the same pattern:
\begin{itemize}
\item \kbd{from LibRobot import *} first. It opens the window, puts the
  robot in the centre, pointing north, pen down, in black, with a grid of
  12 points (the settings of the book's programs). \kbd{from LeRobot import *}
  gives the original interface used by \kbd{my\_drawing.py}: a grid of 30
  points, and the program must start with \kbd{init\_pict(); gc()}.
\item The robot starts with its pen down, so every move draws: call
  \kbd{pu()} to move without drawing, and \kbd{pd()} to draw again.
\item \kbd{st()} last: shows the speed screen and keeps the window open.
  Without it, the bar says \emph{End of the program: press a key or click}
  and the window closes at the next key or click.
\item Commands without arguments take \kbd{()}; commands are separated by
  \kbd{;} or written on separate lines.
\end{itemize}

\subsection{Programs of the book}
The folder \kbd{book} holds the 406 programs of the book \emph{First
Programming Lessons in Python}, one subfolder per chapter:
\begin{lstlisting}[language=bash]
./robot.sh book/20_polar_grid/LimaconOfPascal.py
make run PROG=book/21_recursion/DragonCurve
make check                         # all programs, one after the other
make check DIR=02_interaction      # one chapter
make check FROM=FriezeOfTriangles  # resume at a program
\end{lstlisting}
With \kbd{make check}, answer in the robot window when a drawing is
finished: Return = correct (next program), \kbd{r} = run it again, \kbd{q}
or Esc = stop (the command to resume is printed), any other key starts a
remark, ended by Return, noted in \kbd{check.txt}. A program that never
ends by itself (\kbd{Loop1}, \kbd{Loop2}, \kbd{DigitalClock}) is stopped
with Esc. Programs that ask questions show them in the bottom bar and in
the Terminal; answer in either. \kbd{LEROBOT\_MUET=1} turns the sound off.
These programs are the English translations of those of the French
edition, and draw exactly the same pictures; \kbd{tools/renames} gives the
French name of each one.

\subsection{Replaying at another speed}\label{sec:replay}
When the program reaches \kbd{st()}, the bar shows:
\begin{lstlisting}
+: faster   -: slower   speed: N     other key: redraw   Esc: quit
\end{lstlisting}
\begin{itemize}
\item \kbd{+} / \kbd{-} --- change the speed from 0 to 10 (default 10, the
  fastest). Only the number changes.
\item Any other key, for example the space bar --- clear the window and
  redraw at the chosen speed; the speed screen comes back when it is done.
\item Esc --- close the window.
\end{itemize}
At speed 1 each command takes about 0.1\,s. Speed 0 is step-by-step mode:
the robot does one command each time you press a key. While the robot
draws, Esc (or \kbd{\S} or \kbd{\#}) stops the program.

\subsection{Command reference}
One step is one grid square. The underlined letters give the mnemonic.
Commands with two numbers take them as two arguments: \kbd{moveTo(x, y)}.

\begin{longtable}{@{}p{0.42\linewidth} p{0.53\linewidth}@{}}
\toprule
\textbf{Command} & \textbf{Meaning} \\
\midrule
\endhead
\bottomrule
\endfoot
\kbd{fd()} & Move \underline{f}orwar\underline{d} one step \\
\kbd{fdn(n)} & Move \underline{f}orwar\underline{d} $n$ steps (backwards if $n < 0$) \\
\kbd{cr()} / \kbd{cl()} & \underline{C}urve \underline{r}ight / \underline{l}eft (draws a quarter circle) \\
\kbd{etr()} / \kbd{etl()} & \underline{E}ighth \underline{t}urn \underline{r}ight / \underline{l}eft ($45^\circ$) \\
\kbd{qtr()} / \kbd{qtl()} & \underline{Q}uarter \underline{t}urn \underline{r}ight / \underline{l}eft ($90^\circ$) \\
\kbd{tetr()} / \kbd{tetl()} & \underline{T}hree \underline{e}ighths of a \underline{t}urn \underline{r}ight / \underline{l}eft ($135^\circ$) \\
\kbd{ht()} & \underline{H}alf \underline{t}urn ($180^\circ$) \\
\kbd{fn()} \kbd{fne()} \kbd{fe()} \kbd{fse()} \kbd{fs()} \kbd{fsw()} \kbd{fw()} \kbd{fnw()} & \underline{F}ace \underline{n}orth, NE, E, SE, S, SW, W, NW \\
\kbd{face(n)} & Face direction $n$ (\kbd{north} = 0, \kbd{northEast} = 1 \dots\ \kbd{northWest} = 7) \\
\kbd{rot(a)} & \underline{Rot}ate by $a$ degrees (anticlockwise), keeping the step length \\
\kbd{pd()} / \kbd{pu()} & \underline{P}en \underline{d}own / \underline{u}p \\
\kbd{dp()} / \kbd{dc()} & \underline{D}raw a \underline{p}oint (a square of side \kbd{pw}) / a small \underline{c}ross \\
\kbd{pc(n)} & \underline{P}en \underline{c}olour $n$ (\kbd{white} = 0, \kbd{yellow}, \kbd{blue}, \kbd{violet}, \kbd{green}, \kbd{red}, \kbd{indigo}, \kbd{black} = 7) \\
\kbd{pw(n)} & \underline{P}en \underline{w}idth $n$ (a square pen, as on the Macintosh) \\
\kbd{bg(n)} & Clear the window and fill it with \underline{b}ack\underline{g}round colour $n$ \\
\kbd{paint()} & Paint the closed area behind the robot with the pen colour \\
\kbd{translate(dx, dy)} & Move $dx$ steps across and $dy$ steps up, without drawing \\
\kbd{drawTo(x, y)} & Go to window pixel $(x, y)$, drawing if the pen is down \\
\kbd{moveTo(x, y)} & Go to grid point $(x, y)$, counted from the centre \\
\kbd{fdl(l)} & Move \underline{f}orwar\underline{d} a \underline{l}ength of $l$ points, the grid following the robot \\
\kbd{gc()} & \underline{G}o to the \underline{c}entre \\
\kbd{gt()} \kbd{gb()} \kbd{gr()} \kbd{gl()} & \underline{G}o to the \underline{t}op, \underline{b}ottom, \underline{r}ight or \underline{l}eft edge \\
\kbd{dg()} & \underline{D}raw the \underline{g}rid points \\
\kbd{gs(x)} / \kbd{gsX(x)} / \kbd{gsY(x)} & Set the \underline{g}rid \underline{s}ize (both directions / across / up) \\
\kbd{gsRT(r, t)} & Grid step of length $r$ at angle $t$ (polar grid) \\
\kbd{showRobot()} / \kbd{hideRobot()} & Show / hide the robot \\
\kbd{allowExits()} / \kbd{forbidExits()} & Allow / forbid moving past the window edge \\
\kbd{sp(n)} & Set the \underline{sp}eed (0 to 10) \\
\kbd{ss()} / \kbd{ex()} & \underline{S}tep by \underline{s}tep / normal \underline{ex}ecution (full speed) \\
\kbd{message("text")} & Write a message in the bottom bar \\
\kbd{label("text")} & Write a text at the robot's position \\
\kbd{readInt()} / \kbd{readChar()} & Read an integer / a key \\
\kbd{write(\dots)} / \kbd{writeln(\dots)} & Write in the Terminal \\
\kbd{st()} & End (\underline{st}op): show the speed screen \\
\end{longtable}

Test functions such as \kbd{tb()} (\underline{t}est \underline{b}order:
would the next step leave the window?), \kbd{tn()} \dots\ \kbd{tnw()}
(\underline{t}est whether the robot faces \underline{n}orth \dots) and the
values \kbd{valX()}, \kbd{valGsX()}, \kbd{valPc()}, \kbd{valOrientation()},
\dots\ are also available; so are the music commands (\kbd{note(C, quarter)},
\kbd{rest(half)}, \kbd{metronome(n)}, \kbd{dynamics(n)}), the texts
(\kbd{textFont}, \kbd{textSize}, \kbd{intToString}, \dots) and the time
(\kbd{clock()}, \kbd{date()}, \kbd{delay(t)}). The complete list, with the
French name of each command, is in the \kbd{README.md} file of
\kbd{LeRobot\_Python\_EN} (table \kbd{NOMS} of \kbd{noms.py}).

%======================================================================
\section{Troubleshooting}

\begin{longtable}{@{}p{0.27\linewidth} p{0.33\linewidth} p{0.34\linewidth}@{}}
\toprule
\textbf{Symptom} & \textbf{Cause} & \textbf{Fix} \\
\midrule
\endhead
\bottomrule
\endfoot
The program hangs, no window & No Python with Tk 8.6 on this Mac & Install Python from python.org (installation manual) \\
\addlinespace
\emph{No module named 'LibRobot'} & The program is not in \kbd{LeRobot\_Python\_EN} and was started with \kbd{python3} & Use \kbd{./robot.sh Program.py} or \kbd{LeRobotPythonEN.app} \\
\addlinespace
\emph{NameError: name 'av' is not defined} & A program written for the French edition & Run it with \kbd{LeRobot\_Python\_FR}, or replace the French names (\kbd{tools/fr2en.py} does it) \\
\addlinespace
\emph{Error line N of Prog.py: the robot would leave the frame (fd)} & The robot was asked to move past the window edge & Change the program, or call \kbd{allowExits()} \\
\addlinespace
\kbd{** EXIT FORBIDDEN ** translate} in the terminal & In Draw, a move past the window edge & Only a warning \\
\addlinespace
\emph{Error line N of Prog.py: endless recursion (stack full)} & A procedure that calls itself forever & Check the stopping case of the recursion \\
\addlinespace
The program never ends & An endless loop (\kbd{while True}) & Press Esc in the robot window \\
\addlinespace
Keys do nothing in the window & The window does not have keyboard focus & Click in the window \\
\addlinespace
A saved drawing was lost & Saving from Draw always overwrites \kbd{my\_drawing.py} & Copy \kbd{my\_drawing.py} under a new name before saving again \\
\addlinespace
Text font cannot be changed & \kbd{textFont} and \kbd{textStyle} have no effect & None; texts use Helvetica 13 \\
\end{longtable}

\end{document}
